\chapter{Conclusions and Outlook}\label{ch:m-conclusion}

This thesis developed a bottom-up method for locating favorable distributed-GPU communication paths and testing which observations survive composition and complete solver execution. A unified benchmark suite measured point-to-point, halo, allreduce, and all-to-all patterns, then MoE and CG-inspired application microbenchmarks, through CUDA-aware and SYCL-aware MPI, NCCL, oneCCL, NVSHMEM, and OSHMPI, including persistent kernel-invoked NVSHMEM neighbor kernels. The software contributions extend oneCCL with grouped NCCL point-to-point handling, an OSHMPI backend for the benchmarked operation set, and a validated NVSHMEM backend foundation, and extend aCG with an OSHMPI communicator and a SYCL implementation. Together, these five contributions connect interface engineering and controlled communication measurements to two large sparse matrices through 32 A100 GPUs.

Chapter~\ref{ch:m-discussion} answers the three research questions against that evidence. Read together, they give one conclusion: on this platform device initiation is a conditional tool rather than a general improvement. It repays its cost where communication is repeated, regular, and can remain inside persistent device work, which among the measured patterns is the batched multi-node halo and, in the solver, the monolithic variant at larger process counts. Elsewhere a host-driven library remains competitive or better, because the factors that decide the outcome are memory placement, completion semantics, and progress rather than the location of the call. The benchmark evidence transfers to the application only where the operation and its dependencies are genuinely the same; where the application adds irregular peers or overlap, the benchmark bounds the answer rather than predicting it. These conclusions describe configured paths on proxy-assisted networking, not an intrinsic ordering of libraries or an initiation-only speedup.

The preliminary NCCL~2.31.2 study in Appendix~\ref{app:nccl-preliminary} adds a candidate rather than a revised application result. Symmetric buffers improve the supplied small-allreduce medians relative to NCCL~2.18.5, but the advantage reverses at some intermediate sizes and does not extend to small all-to-all. This reinforces the need to record version, memory handling, and selected operation path before interpreting a library upgrade as evidence for GPU initiation.

The resulting recommendation is to select configured paths for a specific application signature. For GPU initiation, developers should first identify repeated communication that can remain inside persistent device work, then verify whether the target transport is GPU-posted or CPU-proxied. They should benchmark the matching pattern and completion boundary, retain several plausible host- and GPU-initiated stacks, and validate them with a composite or complete application whose synchronization and overlap semantics match production execution.

\section{Future Work}

\subsection{Priority Experiments}\label{sec:priority-experiments}

The reconciled analysis identifies the following priorities before broadening the backend matrix:
\begin{enumerate}
\item \emph{Does staged OSHMPI retain its advantage through GPU-visible completion?} Validate the implemented device-to-host/reduction/host-to-device path in both \texttt{cg\_step} and aCG. Compare the solver's MPI and staged OpenSHMEM scalar options with its halo held fixed, using one-double sums and matched runtime settings.
\item \emph{Does an application-like halo improve prediction?} Replay the directed communication graphs exported from the fixed sparse partitions. Measure submission, completion, and exposed wait with and without independent computation. The replay inputs are prepared; replay performance remains unmeasured.
\item \emph{What portion of the persistent-halo result belongs to invocation?} Compare host-called, host-on-stream, nonpersistent device-invoked, and persistent device-invoked NVSHMEM with message layout, memory, transport, signaling, batch length, and completion held fixed.
\item \emph{Does the favorable regime survive application-like batching?} Sweep batch length and add an interleaved halo--compute--reduction schedule rather than relying only on 100 consecutive exchanges.
\item \emph{Which application change explains the device-solver trend?} Match host- and device-NVSHMEM solver organization and collective dispatch, with separate controls for persistent computation and fusion.
\end{enumerate}

The complete priority report, including CUDA--SYCL correctness repair, matched NVSHMEM dispatch, GPU-posted networking, and conditional NCCL and all-to-all studies, is available in \path{experiments/missing-experiments-report.md}.

\subsection{Broader Directions}

The first priority is to improve experimental provenance. Future campaigns should store source revisions, dependency commits, module versions, node names, GPU and NIC identifiers, rank placement, environment variables, and launcher commands in every result artifact. All cells should use the same number of independent allocations, and reports should include job-level distributions and confidence intervals rather than relying primarily on fastest trials.

The communication stacks also require matched sensitivity experiments. Repeating host NVSHMEM with both collective-dispatch settings would make its scalar comparison reproducible at the same level as MPI and NCCL. Forcing UCC alternatives around the 4\,KiB boundary would test the source-based algorithm explanation. GPU/NIC affinity and progress should be recorded independently; the existing threshold and rail controls constrain their tested configurations rather than closing those questions for every application. The anomalous single-rank allreduce control and the mixed-campaign all-to-all archive also warrant targeted confirmation. MoE itself should be repeated with its own UCC control, since its current setting was inherited from the all-to-all experiment rather than measured on the variable-count path.

The halo replay should preserve the directed per-peer counts rather than approximate them by their averages. The retained graphs reach 17 outgoing peers on an individual rank and show substantial per-rank volume imbalance. Allocation-level distributions can test whether the MPI slow regimes recur under that structure. Supporting kernel controls should compare \texttt{sycl::reduction} with the CUDA dot implementation and ordinary with symmetric device allocations, keeping launch shape and runtime initialization fixed. These controls address the reported compute differences without changing allocator, operand residency, and collective configuration together.

For OSHMPI, the next implementation step is a genuinely nonblocking halo protocol using \texttt{shmem\_putmem\_nbi} in the begin phase and a deferred \texttt{shmem\_quiet} in the end phase. Instrumentation should separate CPU issue time, network progress, and exposed wait before and after the change. Neighbor-only completion should be revisited only after these components are measured. The oneCCL/OSHMPI backend should also evaluate its direct-buffer path against conforming symmetric staging and add robust support for communicator subteams where the OpenSHMEM team lifecycle permits it. The oneCCL/NVSHMEM prototype must first complete dependency import, returned-event semantics, same-team stream ordering, and the initial barrier, allgather, allreduce, all-to-all, and broadcast mappings. Only after those host-enqueued operations are correct and measured should a kernel-initiated path be added and compared against direct NVSHMEM and the existing oneCCL backends.

Direct GPU-posted networking remains an important extension, but the evaluated NVSHMEM~2.11.0 configuration cannot establish its performance. A site-supported experiment needs release-specific prerequisites and a verified active transport, not only an environment-variable request. Host-on-stream, device-side, and monolithic solver variants should use matched collective settings. For NCCL, the stalled run's diagnostics and exact GDAKI prerequisites must be recovered before attributing failure to the NVSHMEM driver limitation. Any user device-API experiment must explicitly establish invocation mode and distinguish intra-node LSA from network GIN behavior.

Finally, the evaluation should extend beyond Leonardo. Repeating the methodology on AMD and Intel accelerators with RCCL, ROC\_SHMEM, Intel SHMEM, and appropriate MPI implementations would test whether the observed decision boundaries generalize. Additional sparse matrices, pipelined CG, larger GPU counts, and a real redistribution-heavy application for the all-to-all benchmark would broaden application validity. A longer-term goal is to turn the measured decision surface into a tool that accepts an application's communication signature and recommends a small set of backend configurations for confirmatory testing.
